% -----------------------------------------------------------------------------
% This file is automatically generated.
% Do not edit manually.
% Source: notebooks/support/probability-distributions/support.Rmd
% -----------------------------------------------------------------------------



\noindent Note that the syntax is more elaborate than what we have done so far. In \ttblue{R} we can choose between passing the values of the arguments only and passing values specifically assigned to the arguments that the function uses.

\noindent There are two advantages to using argument names and values. First, the command is easier to interpret by a reviewer. Second, we can pass the arguments in any order. The following commands are therefore equivalent:


\par\addvspace{\topsep}
\begingroup
\fvset{listparameters={%
  \setlength{\topsep}{0pt}%
  \setlength{\partopsep}{0pt}%
  \setlength{\parsep}{0pt}%
  \setlength{\itemsep}{0pt}%
}}
\begin{Verbatim}[breaklines=true,formatcom=\color{ada_blue}]
norm.cdf(1012, 1030, 115.26 / sqrt(200))
norm.cdf(x= 1012, loc= 1030, scale= 115.26 / sqrt(200))
norm.cdf(scale= 115.26 / sqrt(200), loc= 1030, x= 1012)
\end{Verbatim}
\begin{Verbatim}[breaklines=true,formatcom=\color{ada_light_blue}]
[1] 0.01360269

[1] 0.01360269

[1] 0.01360269
\end{Verbatim}
\endgroup
\par\addvspace{\topsep}
